\begin{abstract}
This paper introduces the \emph{DRM Transformer}, a decoder-only Transformer
whose internal attention space is equipped with a \emph{Directional Relational
Manifold} (DRM) geometry. Instead of relying only on dot products in a fixed
Euclidean space, the model projects queries and keys into a low-dimensional
relational manifold, learns a point-dependent metric tensor, applies
gravitational deformations through token mass, and modulates distances through
a Lorentz-like relativistic factor. The implementation preserves the
autoregressive interface of a language model, while replacing attention scores
with learned geodesic proximity. We also introduce geometric losses to prevent
manifold collapse and to investigate toroidal signatures in dimension four. In
a reproducible baseline with approximately 3.5 million parameters trained on
English Wikipedia, topological analysis with Voronoi cells, LTSA, and
persistent homology suggests that training moves the DRM manifold from an
initial high-genus unstable topology toward a homological signature compatible
with a near-toroidal closed surface. In two out of three trained seeds, the
median topology matches $H_1=2$, $H_2=1$, the expected signature of $T^2$; all
trained seeds show a \emph{Near T2} signal, while random initialization does
not. The strict criterion $f_{T^2}\geq 0.60$ has not yet been reached; therefore
we describe the result as preliminary evidence of a near-toroidal signature,
not as final validation of a stable torus. We discuss the mathematical
formulation, architecture, preliminary results, and limitations of the method.
\end{abstract}
